\section{Experimental Results}

\subsection{Main Accuracy Comparison}
Table~\ref{tab:main} compares SRPA-YOLO with general-purpose compact detectors and UAV-oriented detectors on the deduplicated DUT Anti-UAV test set. SRPA-YOLO achieves 95.850\% Precision, 92.201\% Recall, 95.185\% AP@0.50, 67.071\% \mapall{}, and 77.329\% AP@0.75. YOLOv8n is the direct structural baseline because SRPA-YOLO retains its one-stage training and deployment framework. Relative to YOLOv8n, SRPA-YOLO improves the five metrics by 3.088, 8.242, 5.057, 6.895, and 9.750 percentage points, respectively, while reducing deployed parameters by 0.842 M. Under the common evaluation protocol, the concurrent Precision and Recall gains make an explanation based solely on a shifted confidence operating point less plausible.

\begin{table*}[!t]
\caption{Comparison on 2199 deduplicated DUT Anti-UAV test images. Detection metrics are reported in percent; best values are bold\label{tab:main}}
\centering\srpatablesetup
\begin{tabular}{lrrrrrr}
\toprule
Method & \shortstack{Params\\(M)} & \shortstack{Precision\\(\%)} & \shortstack{Recall\\(\%)} & \shortstack{AP@0.50\\(\%)} & \shortstack{\mapall{}\\(\%)} & \shortstack{AP@0.75\\(\%)}\\
\midrule
YOLO11n & 2.582 & 92.682 & 83.868 & 90.100 & 59.280 & 66.383\\
YOLOv8n & 3.006 & 92.762 & 83.959 & 90.128 & 60.176 & 67.579\\
EDGS-YOLOv8 & 3.091 & 94.266 & 83.244 & 90.428 & 60.347 & 66.415\\
DRBD-YOLOv8 & 3.047 & 92.658 & 83.601 & 90.203 & 60.689 & 67.459\\
YOLOv8n-P2 & 2.921 & \textbf{96.114} & 90.375 & 93.977 & 65.407 & 76.290\\
RLRD-YOLO & 2.926 & 95.781 & 90.107 & 94.467 & 65.436 & 74.757\\
\textbf{SRPA-YOLO} & \textbf{2.164} & 95.850 & \textbf{92.201} & \textbf{95.185} & \textbf{67.071} & \textbf{77.329}\\
\addlinespace
$\Delta$ SRPA vs. YOLOv8n & $-0.842$ & $+3.088$ & $+8.242$ & $+5.057$ & $+6.895$ & $+9.750$\\
\bottomrule
\end{tabular}
\end{table*}


Figure~\ref{fig:comparison} summarizes the accuracy and compactness comparison across the evaluated detectors.

\begin{figure*}[!t]
\centering
\includegraphics[width=0.86\textwidth]{Fig5.pdf}
\caption{Accuracy and compactness comparison on the deduplicated DUT Anti-UAV test set. All detection metrics are reported in percent.\label{fig:comparison}}
\end{figure*}

YOLOv8n-P2 controls for the effect of adding a high-resolution prediction level. Although its Precision is 0.264 percentage points higher, SRPA-YOLO improves Recall, AP@0.50, \mapall{}, and AP@0.75 by 1.826, 1.208, 1.664, and 1.039 percentage points, respectively. This comparison shows that task-feature allocation inside the detection head contributes additional gains beyond high-resolution scale extension.

Figure~\ref{fig:complexity} relates deployed parameters to the primary accuracy metric. SRPA-YOLO occupies the upper-left region, combining the highest measured \mapall{} with the smallest parameter count among the compared detectors. This operating point is consistent with the design objective of assigning a compact shared dictionary and limiting additional capacity to the classification task.

\begin{figure}[!t]
\centering
\includegraphics[width=0.90\columnwidth]{Fig6.pdf}
\caption{Model-complexity trade-off. The horizontal axis is deployed parameter count and the vertical axis is \mapall{}.\label{fig:complexity}}
\end{figure}


\subsection{Shared/Private-Rank and Mechanism Ablations}
Table~\ref{tab:rank} isolates the effect of shared and private ranks on the validation set. Increasing the shared rank from 48 to 64 or 80 does not lead to a monotonic increase in \mapall{}, indicating that uniform width expansion is not sufficient to determine the optimum. The 48+16 allocation reaches the highest validation \mapall{} (60.874\%), whereas 64+16 produces the highest AP@0.75 (69.986\%), AP@0.50 (91.850\%), and Precision (96.615\%). According to Equation~\eqref{eq:selection}, both allocations lie within 0.3 percentage points of the highest \mapall{}, and 64+16 is selected because it provides the strongest high-IoU localization.

\begin{table*}[!t]
\caption{Validation-selected shared/private rank ablation. Detection metrics are percentages; $r_s+r_p$ denotes shared rank plus classification-private rank\label{tab:rank}}
\centering\srpatablesetup
\begin{tabular}{rrrrrrr}
\toprule
$r_s$ & $r_p$ & \shortstack{Params\\(M)} & \shortstack{Precision\\(\%)} & \shortstack{Recall\\(\%)} & \shortstack{AP@0.50\\(\%)} & \shortstack{\mapall{} /\\AP@0.75 (\%)}\\
\midrule
48 & 0  & 2.093 & 96.089 & 87.182 & 91.667 & 60.508 / 69.337\\
64 & 0  & 2.128 & 96.378 & 86.570 & 91.703 & 60.454 / 69.754\\
80 & 0  & 2.164 & 96.073 & \textbf{87.743} & 91.569 & 60.513 / 69.262\\
48 & 16 & 2.128 & 96.135 & 87.310 & 91.798 & \textbf{60.874} / 69.618\\
\textbf{64} & \textbf{16} & \textbf{2.164} & \textbf{96.615} & 86.570 & \textbf{91.850} & 60.646 / \textbf{69.986}\\
\bottomrule
\end{tabular}
\end{table*}


Table~\ref{tab:mechanism} traces cumulative changes from a nonlinear head to the final SRPA configuration. Adding dense DFL distillation is associated with recovery from the accuracy loss of head compression, and the subsequent confidence-focused term adds 0.384 \mapall{} points. Under the cumulative protocol, transferred initialization coincides with increases of 2.660 Recall points and 0.678 \mapall{} points; increasing the shared representation to 64 channels is followed by gains of 0.531 \mapall{} points and 1.674 AP@0.75 points. The final 16-channel classification-private representation reaches 67.071\% \mapall{} and 77.329\% AP@0.75. Each row retains the preceding components and therefore reports the performance of the corresponding cumulative configuration. The shared-80 control reaches a similar \mapall{} of 67.121\%, whereas the 64+16 allocation additionally enforces the zero box Jacobian in Equation~\eqref{eq:jacobian} and is selected by the predefined validation localization criterion. Together with the rank ablation, the results support task-constrained capacity allocation instead of unrestricted head widening.

\begin{table*}[!t]
\caption{Mechanism ablation on the deduplicated DUT Anti-UAV test set. Detection metrics are percentages\label{tab:mechanism}}
\centering\srpatablesetup
\begin{tabular}{p{0.27\textwidth}rrrrrr}
\toprule
Configuration & \shortstack{Params\\(M)} & \shortstack{Precision\\(\%)} & \shortstack{Recall\\(\%)} & \shortstack{AP@0.50\\(\%)} & \shortstack{\mapall{}\\(\%)} & \shortstack{AP@0.75\\(\%)}\\
\midrule
Full nonlinear head & 2.603 & 95.922 & 90.954 & 94.736 & 66.046 & 75.503\\
Shared-48 + dense DFL KD & 2.093 & 96.152 & 90.954 & 94.836 & 65.363 & 74.929\\
$+$ focused localization KD & 2.093 & 97.084 & 90.285 & 95.018 & 65.747 & 75.573\\
$+$ transferred initialization & 2.093 & 96.484 & 92.945 & \textbf{95.493} & 66.425 & 75.472\\
$+$ shared rank 64 & 2.128 & 96.122 & 92.157 & 95.366 & 66.956 & 77.146\\
$+$ private rank 16 & 2.164 & 95.850 & 92.201 & 95.185 & \textbf{67.071} & \textbf{77.329}\\
\bottomrule
\end{tabular}
\end{table*}



\subsection{Localization and Scale Characteristics}
Figure~\ref{fig:scaleiou} compares SRPA-YOLO directly with YOLOv8n over IoU thresholds from 0.50 to 0.95. SRPA-YOLO has higher AP at every threshold. The margin is 5.057 percentage points at IoU 0.50, increases to 9.750 points at IoU 0.75, and remains 8.803 points at IoU 0.85. The persistence of the margin under stricter overlap criteria indicates that the measured improvement includes a localization component and is not confined to class confidence at IoU 0.50.

\begin{figure*}[!t]
\centering
\includegraphics[width=0.84\textwidth]{Fig7.pdf}
\caption{Localization comparison between SRPA-YOLO and YOLOv8n. Left: AP at individual IoU thresholds; right: aggregate detection metrics.\label{fig:scaleiou}}
\end{figure*}

The largest aggregate improvement over YOLOv8n is observed in Recall (+8.242 points), accompanied by gains in both AP@0.50 and AP@0.75. In long-range UAV monitoring, improved Recall directly reduces missed targets, while the concurrent Precision increase limits additional false alarms. The P2 prediction level provides spatial support for tiny UAVs, while the SRPA head keeps classification-private refinement outside the box-regression input. This structural interpretation is consistent with the observed gains at both low and high IoU thresholds.


\subsection{Qualitative Detection Analysis}
Figure~\ref{fig:detectioncomparison} compares YOLOv8n and SRPA-YOLO on the same tiny-UAV instances at a confidence threshold of 0.25. A detection is regarded as correct when its IoU with the ground-truth box is at least 0.50. In all three selected examples, YOLOv8n produces no matched prediction, whereas SRPA-YOLO recovers the target. These representative cases visually corroborate the quantitative improvements reported on the test set. The enlarged regions show the spatial alignment of the recovered predictions for targets represented by only a few pixels.

\begin{figure*}[!t]
\centering
\includegraphics[width=0.84\textwidth]{Fig8.pdf}
\caption{Matched-image detection comparison on the deduplicated DUT Anti-UAV test set. Yellow dashed boxes denote ground truth, red solid boxes denote SRPA-YOLO predictions, and ``Missed'' indicates that YOLOv8n has no prediction satisfying confidence $\geq0.25$ and IoU $\geq0.50$.\label{fig:detectioncomparison}}
\end{figure*}

Figure~\ref{fig:tinydetections} further presents successful SRPA-YOLO detections under sky, building, vegetation, and low-contrast backgrounds. The selected targets have projected longest-edge lengths of 5.7--6.3 pixels after 640-pixel letterboxing. Despite the limited target support and substantial surrounding background, the final validation-selected checkpoint produces localized responses in each selected scene. These cases illustrate the model's small-target detection behavior and are consistent with the preceding test-set metrics.

\begin{figure*}[!t]
\centering
\includegraphics[width=0.74\textwidth]{Fig9.pdf}
\caption{Successful SRPA-YOLO detections of tiny UAVs. Each right-hand view enlarges the target region from the corresponding full image; red boxes and labels are predictions from the final validation-selected checkpoint.\label{fig:tinydetections}}
\end{figure*}


\subsection{Matched-Seed Stability}
Table~\ref{tab:seeds} summarizes the results for Stage-II seeds 40, 41, and 42. All experiments use the same Stage-I shared basis and independently initialize the classification-private parameters. SRPA-YOLO yields standard deviations of 0.112, 0.022, and 0.082 percentage points for Recall, AP@0.50, and \mapall{}, respectively, indicating low variation across these private-parameter initializations. Relative to RLRD-YOLO, the mean paired differences in Recall, \mapall{}, and AP@0.75 are +2.709, +1.881, and +2.150 percentage points, respectively. Under the matched protocol, the observed improvement is consistent across the three Stage-II seeds.

\begin{table}[H]
\caption{Matched Stage-II seeds on the deduplicated DUT Anti-UAV test set (percent, mean $\pm$ sample standard deviation, $n=3$). All runs use the same Stage-I shared basis and independently initialize only the private parameters\label{tab:seeds}}
\centering\srpatablesetup
\begin{tabular}{@{}lrrr@{}}
\toprule
Metric (\%) & RLRD-YOLO & SRPA-YOLO & Paired diff.\\
\midrule
Precision & $95.964\pm0.766$ & $\mathbf{96.167\pm0.031}$ & $+0.204\pm0.736$\\
Recall & $89.463\pm0.785$ & $\mathbf{92.172\pm0.112}$ & $+2.709\pm0.773$\\
AP@0.50 & $94.661\pm0.106$ & $\mathbf{95.304\pm0.022}$ & $+0.644\pm0.109$\\
\mapall{} & $65.199\pm0.639$ & $\mathbf{67.080\pm0.082}$ & $+1.881\pm0.721$\\
AP@0.75 & $75.196\pm0.564$ & $\mathbf{77.345\pm0.233}$ & $+2.150\pm0.332$\\
\bottomrule
\end{tabular}
\end{table}
\vspace{-4pt}



\subsection{Latency and Deployment Complexity}
The training graph contains approximately 2.18 M parameters and requires 8.5 GFLOPs at 640 pixels. Structural fusion removes training-only state and produces a 2.164 M-parameter checkpoint. Table~\ref{tab:latency} reports same-process paired measurements, in which each detector is evaluated together with a paired YOLOv8n control; Table~\ref{tab:cleanlat} gives the separate low-load final-acceptance measurements.

\begin{table}[!t]
\caption{Same-process paired latency. Each row reports a candidate and its paired YOLOv8n control; batch 1, 640 pixels, FP16, fused forward+NMS, 80 warm-ups, and 300 iterations\label{tab:latency}}
\centering\srpatablesetup
\begin{tabular}{@{}lrrrr@{}}
\toprule
Method & Round & ms & FPS & Control ms\\
\midrule
YOLO11n & 1 & 21.6357 & 46.22 & 16.8632\\
 & 2 & 19.1505 & 52.22 & 15.0139\\
YOLOv8n-P2 & 1 & 19.1953 & 52.10 & 15.9227\\
 & 2 & 16.9629 & 58.95 & 14.1962\\
EDGS* & 1 & 19.5146 & 51.24 & 14.6088\\
 & 2 & 22.0043 & 45.45 & 16.3604\\
DRBD* & 1 & 29.0158 & 34.46 & 14.5733\\
 & 2 & 31.7115 & 31.53 & 16.8584\\
RLRD-YOLO & 1 & 22.5644 & 44.32 & 16.0771\\
 & 2 & 22.7773 & 43.90 & 16.2647\\
\textbf{SRPA-YOLO} & \textbf{1} & \textbf{16.2469} & \textbf{61.55} & 16.4922\\
 & \textbf{2} & \textbf{15.0936} & \textbf{66.25} & 15.5909\\
\bottomrule
\end{tabular}
\end{table}

\begin{table}[!t]
\caption{Separate low-load final-acceptance measurements for SRPA-YOLO and its paired YOLOv8n control under the same batch-1, 640-pixel, FP16, fused forward+NMS protocol\label{tab:cleanlat}}
\centering\srpatablesetup
\begin{tabular}{@{}lrrrr@{}}
\toprule
Round & \shortstack{SRPA\\ms} & \shortstack{SRPA\\FPS} & \shortstack{YOLOv8n\\ms} & \shortstack{YOLOv8n\\FPS}\\
\midrule
1 & \textbf{9.7435} & \textbf{102.63} & 9.8859 & 101.15\\
2 & \textbf{9.8574} & \textbf{101.45} & 9.9893 & 100.11\\
\bottomrule
\end{tabular}
\end{table}


In the separate low-load acceptance measurements, SRPA-YOLO requires 9.7435 and 9.8574 ms (102.63 and 101.45 FPS), whereas the paired YOLOv8n control requires 9.8859 and 9.9893 ms (101.15 and 100.11 FPS). The two-round means are 9.8005 ms and 9.9376 ms, respectively. Under the reported protocol, SRPA-YOLO achieves comparable frame time while reducing the deployed parameter count by 28.0\%.

Figure~\ref{fig:speedtradeoff} combines the reported two-round mean latency with the primary accuracy metric. Among the measured models and under this hardware protocol, SRPA-YOLO occupies the favorable accuracy--speed boundary: models with comparable accuracy have higher measured latency, whereas faster models do not reach the same \mapall{}.

\begin{figure}[!t]
\centering
\includegraphics[width=0.72\columnwidth]{Fig10.pdf}
\caption{Accuracy--speed trade-off based on the reported two-round mean latency and \mapall{}.\label{fig:speedtradeoff}}
\end{figure}

\FloatBarrier
